% ==================== 1.2 国内外研究现状 ====================
\section{国内外研究现状}
\label{sec:research_status}

月面巡视器自主行走涉及月面环境建模、轮-壤相互作用、视觉感知、路径规划等多个学科领域。本节对相关领域的研究现状进行系统梳理和分析。

\subsection{月面环境建模研究现状}
\label{subsec:env_modeling}

月面环境建模是巡视器自主行走的基础，主要包括地形建模、月壤建模和光照温度建模三个方面。

\textbf{（1）月面地形建模}

月面地形建模方法可分为基于实测数据的方法和基于生成模型的方法两大类。基于实测数据的方法主要利用遥感影像、激光测高数据等构建数字高程模型（DEM）。美国地质调查局（USGS）基于LRO（Lunar Reconnaissance Orbiter）数据发布了全球1:100万比例尺的月面地形图\cite{ref7}。国内方面，中科院国家天文台利用嫦娥一号激光高度计数据构建了全月球DEM模型，分辨率优于500米\cite{ref8}。

基于生成模型的方法主要采用分形理论、物理模拟等方法生成逼真的地形数据。Musgrave等\cite{ref9}提出的分形地形生成算法被广泛应用于地外环境模拟。近年来，深度学习方法开始应用于地形生成，如Gao等\cite{ref10}利用生成对抗网络（GAN）生成高分辨率月面地形。

\textbf{（2）月壤建模}

月壤力学特性是影响巡视器行走性能的关键因素。传统方法主要基于半经验的贝克-雷帕乔夫（Bekker-Reece）理论描述轮-壤相互作用\cite{ref11}。Apollo和Luna计划获取的月壤样本为月壤力学参数研究提供了珍贵数据\cite{ref12}。

近年来，离散元方法（DEM）被广泛应用于月壤力学模拟。Smith等\cite{ref13}建立了月壤颗粒的DEM模型，分析了颗粒形状对力学特性的影响。国内方面，北京航空航天大学等机构对模拟月壤进行了大量实验研究\cite{ref14}。

\textbf{（3）光照温度建模}

月面光照和温度环境对巡视器能源供给和热控设计至关重要。Vaniman等\cite{ref15}建立了月面光照条件分析模型。Paige等\cite{ref16}基于Diviner数据绘制了全月球表面温度分布图。国内方面，中科院空间中心对月球极区光照条件进行了深入研究\cite{ref17}。

\subsection{月面巡视器轮-壤相互作用力学研究现状}
\label{subsec:wheel_soil}

轮-壤相互作用力学是巡视器动力学建模的核心。现有研究主要分为理论建模、数值模拟和实验验证三个方面。

\textbf{（1）理论建模}

经典的轮-壤相互作用理论以Bekker\cite{ref11}的地面力学理论为基础，Wong\cite{ref18}对其进行了系统发展和完善。该理论通过沉陷量-压力关系和剪切应力-位移关系描述车轮与土壤的相互作用，建立了车辆-土壤系统力学模型。

针对月面环境特点，研究者们对经典理论进行了改进。Iagnemma等\cite{ref19}考虑了低重力环境对轮-壤相互作用的影响。Ding等\cite{ref20}建立了考虑滑移特性的轮-壤相互作用模型。

\textbf{（2）数值模拟}

有限元法（FEM）和离散元法（DEM）是轮-壤相互作用数值模拟的主要方法。Fervers\cite{ref21}利用FEM分析了车轮沉陷和牵引力特性。Nakashima等\cite{ref22}采用DEM模拟了轮-壤相互作用过程。这些方法可以获取理论模型难以描述的细观力学信息。

\textbf{（3）实验验证}

国内外建立了多个轮-壤相互作用实验平台。NASA Glenn研究中心建立了模拟月球环境的土壤力学实验装置\cite{ref23}。国内哈尔滨工业大学、北京航空航天大学等机构也建立了相关实验系统\cite{ref24}。

\subsection{月面巡视器视觉感知与地形理解研究现状}
\label{subsec:visual_perception}

视觉感知是巡视器自主行走的信息基础，包括环境感知、目标识别和定位导航等方面。

\textbf{（1）环境感知}

月面环境感知主要依赖光学相机和激光雷达等传感器。NASA火星车采用立体相机进行三维地形重建\cite{ref25}。"玉兔号"巡视器搭载的导航相机和全景相机实现了周围环境感知\cite{ref26}。

\textbf{（2）目标识别}

岩石识别是月面目标识别的主要任务。传统方法主要基于视觉特征和几何特征进行识别\cite{ref27}。近年来，深度学习方法在目标识别中展现出优势。Li等\cite{ref28}利用卷积神经网络进行月面岩石识别，取得了较好的效果。

\textbf{（3）定位导航}

视觉SLAM（Simultaneous Localization and Mapping）技术在巡视器定位中得到广泛应用。ORB-SLAM\cite{ref29}、LSD-SLAM\cite{ref30}等代表性算法为巡视器定位提供了技术支撑。

\subsection{月面巡视器路径规划与自主导航研究现状}
\label{subsec:path_planning}

路径规划与自主导航是巡视器自主行走的核心技术。

\textbf{（1）路径规划算法}

传统路径规划算法包括Dijkstra算法、A*算法\cite{ref31}、RRT算法\cite{ref32}等。这些算法在静态环境下取得了良好效果。针对动态环境，研究者提出了D*、D* Lite等增量式规划算法\cite{ref33}。

近年来，基于强化学习的路径规划方法受到关注。深度Q网络（DQN）及其改进算法\cite{ref34,ref35}为复杂环境下的路径规划提供了新思路。

\textbf{（2）自主导航系统}

NASA的自主导航系统已应用于火星探测任务\cite{ref36}。ESA为ExoMars火星车开发了自主导航软件\cite{ref37}。国内方面，中国空间技术研究院研发了月面巡视器自主导航系统\cite{ref38}。

\subsection{数字孪生技术在航天领域的应用现状}
\label{subsec:dt_aerospace}

数字孪生技术在航天领域展现出广阔的应用前景。

\textbf{（1）航天器健康监测}

NASA在"阿耳忒弥斯"计划中采用数字孪生技术对航天器进行健康监测和预测性维护\cite{ref39}。欧洲航天局（ESA）建立了卫星数字孪生系统用于在轨故障诊断\cite{ref40}。

\textbf{（2）任务仿真验证}

NASA建立了航天器组装、测试和发射操作的数字孪生系统\cite{ref41}。波音公司利用数字孪生技术对航空航天系统进行全生命周期管理\cite{ref42}。

\textbf{（3）深空探测应用}

NASA将数字孪生技术应用于火星探测器"毅力号"的操作控制\cite{ref43}。国内方面，中国空间技术研究院开始探索数字孪生在航天器研制中的应用\cite{ref44}。

综上所述，现有研究在各自领域取得了丰富成果，但面向月面巡视器自主行走的数字孪生技术研究尚处于起步阶段，存在以下不足：（1）缺乏面向月面巡视器自主行走的数字孪生理论框架；（2）月面多物理场环境数字孪生建模不够系统；（3）巡视器动力学数字孪生模型的实时性和准确性有待提高；（4）数字孪生环境下的感知、规划与决策方法研究不足。
